% ============================================================
\chapter{Distances entre Diagrammes}
\label{ch:distances}
% ============================================================

Pour que les diagrammes de persistance soient v\'eritablement utiles, il faut pouvoir les \emph{comparer}. Deux nuages de points ont-ils la m\^eme forme topologique~? Un signal bruit\'e a-t-il la m\^eme structure qu'un signal propre~? R\'epondre \`a ces questions exige une notion de distance entre diagrammes. Deux familles dominent la th\'eorie~: la distance bottleneck, qui mesure le pire d\'eplacement n\'ecessaire pour apparier les points de deux diagrammes, et les distances de Wasserstein, qui mesurent le co\^ut total du transport optimal. Le choix entre les deux n'est pas neutre~: la distance bottleneck privil\'egie les structures les plus persistantes (les ``grandes'' barres), tandis que Wasserstein int\`egre l'information de toutes les barres, y compris les petites. Les algorithmes de calcul --- probl\`eme d'affectation hongrois, enchères --- et les propri\'et\'es th\'eoriques de stabilit\'e font l'objet de ce chapitre.

\begin{intuition}
Les diagrammes de persistance\index{diagramme de persistance} r\'esum\'ent
l'information topologique d'un espace filtr\'e. Pour pouvoir comparer deux
espaces, il faut disposer de \textbf{distances} entre diagrammes. Le choix
de la distance d\'etermine quelles diff\'erences topologiques on privil\'egie~:
les structures les plus persistantes (bottleneck\index{distance bottleneck})
ou l'ensemble du spectre (Wasserstein\index{distance de Wasserstein}).
\end{intuition}

% ============================================================
\section{Rappels sur les diagrammes de persistance}
% ============================================================

\begin{definition}[Diagramme de persistance]
Soit $\mathcal{F}$ une filtration d'un complexe simplicial. Le
\textbf{diagramme de persistance}\index{diagramme de persistance}
$\mathrm{Dgm}(\mathcal{F})$ est le multi-ensemble de points
$(b_i, d_i) \in \R^2$ avec $b_i < d_i$, o\`u $b_i$ est la valeur de
naissance et $d_i$ la valeur de mort de la $i$-\`eme classe d'homologie
persistante. On y adjoint la diagonale $\Delta = \{(x,x) : x \in \R\}$
avec multiplicit\'e infinie.
\end{definition}

\begin{remarque}
L'ajout de la diagonale avec multiplicit\'e infinie permet de d\'efinir
des bijections entre diagrammes de tailles diff\'erentes~: un point du
diagramme peut \^etre appari\'e \`a un point de la diagonale, ce qui
revient \`a le consid\'erer comme du bruit.
\end{remarque}

% ============================================================
\section{La distance bottleneck}
\label{sec:bottleneck}
% ============================================================

\begin{definition}[Distance bottleneck]
\label{def:bottleneck}
Soient $D_1$ et $D_2$ deux diagrammes de persistance. La \textbf{distance
bottleneck}\index{distance bottleneck} est~:
\[
  d_\infty(D_1, D_2) = \inf_{\gamma} \sup_{x \in D_1}
    \norm{x - \gamma(x)}_\infty
\]
o\`u l'infimum porte sur toutes les bijections $\gamma : D_1 \to D_2$
(en utilisant la diagonale pour compl\'eter).
\end{definition}

\begin{remarque}
La norme $\norm{\cdot}_\infty$ utilis\'ee ici est la norme du maximum~:
$\norm{(a,b) - (c,d)}_\infty = \max(\abs{a-c}, \abs{b-d})$.
On peut aussi utiliser la distance \`a la diagonale~:
$\norm{(b,d)}_\infty^{\Delta} = \frac{d - b}{2}$.
\end{remarque}

\begin{intuition}
La distance bottleneck mesure le \og{} co\^ut du pire appariement \fg.
Elle est insensible au nombre de petites structures et ne capture que
le d\'ecalage maximal entre les structures les plus importantes.
\end{intuition}

\begin{exemple}
Consid\'erons $D_1 = \{(0, 2), (1, 3)\}$ et $D_2 = \{(0, 2.5), (1, 3.2)\}$.
L'appariement naturel donne~:
\[
  d_\infty(D_1, D_2) = \max\bigl(\norm{(0,2)-(0,2.5)}_\infty,\;
  \norm{(1,3)-(1,3.2)}_\infty\bigr) = \max(0.5, 0.2) = 0.5.
\]
\end{exemple}

\begin{proposition}[M\'etrique]
$d_\infty$ est une m\'etrique sur l'espace des diagrammes de persistance.
\end{proposition}

% ============================================================
\section{Les distances de Wasserstein}
\label{sec:wasserstein}
% ============================================================

\begin{definition}[Distance de Wasserstein]
\label{def:wasserstein}
Pour $p \geq 1$ et $q \geq 1$, la \textbf{distance de Wasserstein}
$W_{p,q}$ entre deux diagrammes $D_1$, $D_2$ est~:
\[
  W_{p,q}(D_1, D_2) = \left(\inf_{\gamma}
    \sum_{x \in D_1} \norm{x - \gamma(x)}_q^p\right)^{1/p}
\]
o\`u $\gamma$ parcourt les bijections $D_1 \to D_2$.
Le cas $q = \infty$ et $p = \infty$ redonne la distance bottleneck.
\end{definition}

\begin{attention}[title=\textbf{Conventions variables}]
La litt\'erature utilise diff\'erentes conventions pour les indices de
la distance de Wasserstein. Certains auteurs \'ecrivent $W_p$ avec
$q = \infty$ implicite, d'autres utilisent $q = p$. V\'erifiez
toujours la convention utilis\'ee dans chaque article.
\end{attention}

\begin{proposition}[Relation bottleneck--Wasserstein]
Pour tout $p \geq 1$~:
\[
  d_\infty(D_1, D_2) = \lim_{p \to \infty} W_{p,\infty}(D_1, D_2).
\]
De plus, $d_\infty(D_1, D_2) \leq W_{p,\infty}(D_1, D_2)$ pour tout $p$.
\end{proposition}

\begin{formules}[title=\textbf{R\'esum\'e des distances}]
\begin{align}
  d_\infty(D_1, D_2) &= \inf_{\gamma} \sup_{x \in D_1}
    \norm{x - \gamma(x)}_\infty \\
  W_{p,q}(D_1, D_2) &= \left(\inf_{\gamma}
    \sum_{x \in D_1} \norm{x - \gamma(x)}_q^p\right)^{1/p}
\end{align}
\end{formules}

% ============================================================
\section{Le th\'eor\`eme de stabilit\'e}
\label{sec:stabilite}
% ============================================================

\begin{theoreme}[Stabilit\'e de la persistance --- Cohen-Steiner, Edelsbrunner, Harer 2007]
\label{thm:stabilite}
\index{th\'eor\`eme de stabilit\'e}
Soient $f, g : X \to \R$ deux fonctions continues sur un espace
topologique triangulable $X$. Alors~:
\[
  d_\infty\bigl(\mathrm{Dgm}(f),\; \mathrm{Dgm}(g)\bigr) \leq
  \norm{f - g}_\infty.
\]
\end{theoreme}

\begin{intuition}
Le th\'eor\`eme de stabilit\'e est le r\'esultat le plus important de la TDA.
Il garantit que de petites perturbations des donn\'ees (au sens $L^\infty$)
ne peuvent produire que de petites perturbations du diagramme de persistance.
C'est ce qui rend la persistance \textbf{robuste au bruit}.
\end{intuition}

\begin{corollaire}
L'application $f \mapsto \mathrm{Dgm}(f)$ est $1$-lipschitzienne de
$\bigl(C(X, \R), \norm{\cdot}_\infty\bigr)$ vers
$\bigl(\mathcal{D}, d_\infty\bigr)$.
\end{corollaire}

\begin{theoreme}[Stabilit\'e en Wasserstein]
Sous des hypoth\`eses de r\'egularit\'e suppl\'ementaires (fonctions de
Morse apprivois\'ees), on a pour tout $p \geq 1$~:
\[
  W_{p,\infty}(\mathrm{Dgm}(f), \mathrm{Dgm}(g)) \leq
  C_p \cdot \norm{f - g}_\infty^{1 - 1/p} \cdot
  \norm{f - g}_{p}^{1/p}
\]
pour une constante $C_p$ d\'ependant de $X$ et $p$.
\end{theoreme}

% ============================================================
\section{La distance d'entrelacement}
\label{sec:interleaving}
% ============================================================

\begin{definition}[Modules de persistance]
\index{module de persistance}
Un \textbf{module de persistance} est un foncteur $\mathbb{V} :
(\R, \leq) \to \mathbf{Vec}_k$, c'est-\`a-dire une famille d'espaces
vectoriels $\{V_t\}_{t \in \R}$ avec des applications lin\'eaires
$v_{s,t} : V_s \to V_t$ pour $s \leq t$ v\'erifiant $v_{t,t} = \mathrm{id}$
et $v_{s,u} = v_{t,u} \circ v_{s,t}$ pour $s \leq t \leq u$.
\end{definition}

\begin{definition}[$\varepsilon$-entrelacement]
\index{entrelacement}
Deux modules de persistance $\mathbb{V}$ et $\mathbb{W}$ sont
\textbf{$\varepsilon$-entrelac\'es} s'il existe des morphismes de
modules d\'ecal\'es $\varphi : \mathbb{V} \to \mathbb{W}[\varepsilon]$
et $\psi : \mathbb{W} \to \mathbb{V}[\varepsilon]$ tels que~:
\[
  \psi[\varepsilon] \circ \varphi = v_{2\varepsilon}, \qquad
  \varphi[\varepsilon] \circ \psi = w_{2\varepsilon}
\]
o\`u $\mathbb{W}[\varepsilon]_t = W_{t+\varepsilon}$ est le module d\'ecal\'e.
\end{definition}

\begin{definition}[Distance d'entrelacement]
La \textbf{distance d'entrelacement}\index{distance d'entrelacement} est~:
\[
  d_I(\mathbb{V}, \mathbb{W}) = \inf\{\varepsilon \geq 0 :
  \mathbb{V} \text{ et } \mathbb{W} \text{ sont }
  \varepsilon\text{-entrelac\'es}\}.
\]
\end{definition}

\begin{theoreme}[Isometrie --- Chazal et al.\ 2009]
\label{thm:isometrie}
\index{th\'eor\`eme d'isom\'etrie}
Pour des modules de persistance \`a param\`etre r\'eel d\'ecomposables
en intervalles~:
\[
  d_I(\mathbb{V}, \mathbb{W}) = d_\infty(\mathrm{Dgm}(\mathbb{V}),
  \mathrm{Dgm}(\mathbb{W})).
\]
\end{theoreme}

\begin{intuition}
Le th\'eor\`eme d'isom\'etrie montre que la distance bottleneck entre
diagrammes co\"\i ncide exactement avec la distance d'entrelacement
entre modules. Cela donne une interpr\'etation \textbf{alg\'ebrique}
de la distance bottleneck et fonde la stabilit\'e sur des bases
cat\'egoriques solides.
\end{intuition}

% ============================================================
\section{Stabilit\'e alg\'ebrique}
\label{sec:stabilite_algebrique}
% ============================================================

\begin{theoreme}[Stabilit\'e alg\'ebrique]
\index{stabilit\'e alg\'ebrique}
Soit $F : \mathbf{Top} \to \mathbf{Vec}_k^{(\R, \leq)}$ un foncteur
d'homologie persistante. Pour toute paire d'espaces filtr\'es
$(X, f)$ et $(Y, g)$ avec un morphisme $\varepsilon$-entrelac\'e
entre les filtrations~:
\[
  d_I(H_*(X, f),\; H_*(Y, g)) \leq \varepsilon.
\]
\end{theoreme}

\begin{remarque}
La stabilit\'e alg\'ebrique g\'en\'eralise le th\'eor\`eme de stabilit\'e
classique au cadre des modules de persistance abstraits. Elle ne
n\'ecessite pas de fonctions continues~: elle s'applique directement
au niveau des filtrations et de l'alg\`ebre des modules.
\end{remarque}

% ============================================================
\section{Calcul pratique des distances}
\label{sec:calcul}
% ============================================================

\begin{proposition}[Complexit\'e]
La distance bottleneck entre deux diagrammes de $n$ et $m$ points
peut \^etre calcul\'ee en $\mathcal{O}(N^{1.5} \log N)$ o\`u
$N = n + m$, par r\'eduction au probl\`eme d'appariement biparti.
La distance de Wasserstein $W_p$ se calcule en $\mathcal{O}(N^3)$
par l'algorithme hongrois.
\end{proposition}

\begin{minted}{python}
import numpy as np
from gudhi.wasserstein import wasserstein_distance
from gudhi.bottleneck import bottleneck_distance

# Deux diagrammes de persistance (naissance, mort)
dgm1 = np.array([[0.0, 2.0], [1.0, 3.0], [2.0, 2.5]])
dgm2 = np.array([[0.0, 2.1], [1.1, 3.2]])

# Distance bottleneck
d_bn = bottleneck_distance(dgm1, dgm2)
print(f"Bottleneck: {d_bn:.4f}")

# Distance de Wasserstein (p=2)
d_w2 = wasserstein_distance(dgm1, dgm2, order=2)
print(f"Wasserstein-2: {d_w2:.4f}")
\end{minted}

% ============================================================
\section{Exercices}
% ============================================================

\begin{exercice}
Calculer \`a la main la distance bottleneck entre
$D_1 = \{(0,1), (0,3), (2,5)\}$ et
$D_2 = \{(0,1.5), (1,4)\}$.
\textit{Indication}~: essayez diff\'erents appariements et n'oubliez
pas la diagonale.
\end{exercice}

\begin{exercice}
Montrer que $d_\infty(D_1, D_2) \leq W_{p,\infty}(D_1, D_2)$
pour tout $p \geq 1$.
\end{exercice}

\begin{exercice}
Soient $f(x) = \sin(x)$ et $g(x) = \sin(x) + 0.1$ sur $[0, 2\pi]$.
Montrer que $d_\infty(\mathrm{Dgm}(f), \mathrm{Dgm}(g)) \leq 0.1$
en utilisant le th\'eor\`eme de stabilit\'e.
\end{exercice}

\begin{exercice}[Distance d'entrelacement]
Soient $\mathbb{V}$ le module $k$ sur $[0, 2]$ et $\mathbb{W}$ le
module $k$ sur $[0.5, 2.5]$. Calculer $d_I(\mathbb{V}, \mathbb{W})$
directement \`a partir de la d\'efinition.
\end{exercice}

\begin{exercice}[Impl\'ementation]
En utilisant \texttt{gudhi}, comparer les distances bottleneck et
Wasserstein-$2$ pour des diagrammes de persistance de deux nuages
de points~: un \'echantillon du cercle $S^1$ et un \'echantillon
du tore $T^2$ plong\'e dans $\R^3$.
\end{exercice}
